% ex: ts=2 sw=2 sts=2 et filetype=tex
% SPDX-License-Identifier: CC-BY-SA-4.0
\documentclass[aspectratio=169]{beamer}

\usetheme{Boadilla}

\setbeamertemplate{navigation symbols}{} % To remove the navigation symbols from the bottom of all slides

\graphicspath{{../img/}{../styles/iteso/}}
% use the command "usepackage" instead of "usecolortheme" to use the path to styles
\usepackage{../styles/iteso/beamercolorthemeiteso}

\usepackage[utf8]{inputenc}
\usepackage{graphicx} % Allows including images

\usepackage{tikz} % the tikz package

\usepackage{smartdiagram}

\usetikzlibrary{shapes.geometric, arrows, shapes.arrows}

\tikzstyle{iniciofin} = [rectangle,
                         rounded corners,
                         minimum width=1.5cm,
                         minimum height=0.5cm,
                         text centered,
                         font=\small,
                         draw=black,
                         fill=orange!20 ]

\tikzstyle{entrada} = [trapezium,
                  trapezium left angle=70,
                  trapezium right angle=110,
                  minimum width=1.5cm,
                  minimum height=0.5cm,
                  text centered,
                  font=\small,
                  draw=black,
                  fill=red!20]

\tikzstyle{salida} = [trapezium,
                  trapezium left angle=70,
                  trapezium right angle=110,
                  minimum width=1.5cm,
                  minimum height=0.5cm,
                  text centered,
                  font=\small,
                  draw=black,
                  fill=green!20]

\tikzstyle{proceso} = [rectangle,
                       minimum width=1.5cm,
                       minimum height=0.5cm,
                       text centered,
                       font=\small,
                       draw=black,
                       fill=yellow!20]

\tikzstyle{decision} = [diamond,
                        minimum width=1.5cm,
                        minimum height=0.5cm,
                        text centered,
                        font=\small,
                        draw=black,
                        fill=blue!20]

\tikzstyle{flecha} = [thick,->,>=stealth]


\title[Conceptos básicos de la programación de computadoras]{Conceptos básicos de la programación de computadoras}

\author{Miguel Bernal Marin} % Your name
\institute[ITESO] % Your institution as it will appear on the bottom of every slide, may be shorthand to save space
{
 ITESO, Universidad \\
 Jesuita de Guadalajara \\% Your institution for the title page
\medskip
\textit{miguel.bernal@iteso.mx} % Your email address
}
\date{
Otoño 2025
} % Date, can be changed to a custom date


%------------------------------------------------------------
%The next block of commands puts the table of contents at the 
%beginning of each section and highlights the current section:
\AtBeginSection[]
{
  \begin{frame}
    \frametitle{Contenido}
    \tableofcontents[currentsection]
  \end{frame}
}
%------------------------------------------------------------

\begin{document}
{ % Portada de la universidad
  \setbeamertemplate{footline}{}
  \usebackgroundtemplate{\includegraphics[width=\paperwidth]{fondo-portada.png}}
  \begin{frame}
  \end{frame}
}

\begin{frame}
    \titlepage
\end{frame}

\usebackgroundtemplate{\includegraphics[width=\paperwidth]{fondo.png}}

% ex: ts=4 sw=4 sts=4 et filetype=tex
% SPDX-License-Identifier: CC-BY-SA-4.0
\begin{frame}
    \frametitle{Contenido}
    \tableofcontents
\end{frame}

\section{Algoritmo y sus características}

\begin{frame}[c]{Título}
    \begin{center}
        Texto
    \end{center}
\end{frame}

\section{Lenguajes de programación y su clasifiación}

\begin{frame}[c]{Título}
    \begin{center}
        Texto
    \end{center}
\end{frame}


{ % Portada de la universidad
  \setbeamertemplate{footline}{}
  \usebackgroundtemplate{\includegraphics[width=\paperwidth]{fondo-portada.png}}
  \begin{frame}
  \end{frame}
}
\end{document}
